\section{Numerical Results}
Fig.~\ref{fig:L_dependence} shows  the   dependence of the norm of quasi TMDWFs on the length $\ell$ of the Wilson-line.  As one can see from this figure, with $\{P^z, b_{\perp}, t\}=\{6\pi/L, 3a, 6a\}$, both the quasi-TMDWF $\phi_\ell(0,b_\perp,P^z,\ell)$   and  the square root of the Wilson loop $Z_E$ decay exponentially with   length $\ell$, but the subtracted quasi-TMDWF is length independent when $\ell\ge 0.4$ fm. Some other cases with larger $P^z$, $b_{\perp}$, and $t$ can be found in the supplemental materials~\cite{supplemental}. Based on this observation, we will use $\ell=7a=0.686$ fm as asymptotic results for all cases in the following calculation.

\begin{figure}[!th]
\begin{center}
\includegraphics[width=0.45\textwidth]{images/03_ff_P3b3.pdf}
\caption{The ratios $C_3(b_\perp,P^z,t_{\rm sep},t)/C_2(0, P^z,0,t_{\rm sep})$ (data points) which converge to the ground state contribution at $t,t_{\rm sep}\rightarrow \infty$ (gray band) as   function of $t_{\rm sep}$ and $t$, with $\{P^z, b_{\perp}\}=\{6\pi/L, 3a\}$. As in this figure, our data in general agree with the predicted fit function (colored bands). }\label{fig:joint_fit}
\end{center}
\end{figure}

We performed a joint fit of the form factor and quasi-TMDWF with the same $P^z$ and  $b_{\perp}$ with the parameterization in Eqs.~\eqref{eq:3pt} and \eqref{eq:2pt}. The ratios $C_3(b_\perp,P^z,t_{\rm sep},t)/C_2(0, P^z,0,t_{\rm sep})$ with different $t_{\rm sep}$ and $t$ for the $\{P^z, b_{\perp}\}=\{6\pi/L, 3a\}$ case are shown in Fig.~\ref{fig:joint_fit}, with ground state contribution (gray band) and the fitted results at finite $t_{2}$ and $t$ (colored bands). In this calculation, the excited state contribution is properly described by the fit with  $\chi^2/{\rm d.o.f.}=0.6$. The details of the joint fit, and also more fit quality checks are shown in the supplemental materials~\cite{supplemental}, with similar fitting quality.

\begin{figure}[!th]
\includegraphics[width=0.45\textwidth]{images/04_sf_perturb_sqrt2.pdf}
\caption{ The intrinsic soft factor as a function of $b_{\perp}$  with  $b_{\perp,0}=a$ as in Eq.~\eqref{eq:S_ratio_z=0}. With different pion momentum $P^z$, the results are consistent with each other. The dashed  curve shows  the result of the 1-loop calculation, see Eq.~(\ref{eq:S_ratio_1loop}),  with the strong coupling constant $\alpha_s(1/b_{\perp})$. {The shaded band corresponds to the scale uncertainty  of $\alpha_s$: $\mu\in [1/\sqrt{2},\sqrt{2}]\times 1/b_\perp$. The systematic uncertainty from the operator mixing has been taken into account.  }}\label{fig:quasi-soft-factor}
\end{figure}

The resulting soft factor as function of $b_{\perp}$ is plotted in Fig.~\ref{fig:quasi-soft-factor}, at $\gamma$= 2.17,  3.06 and 3.98, which corresponds to $P^z=\{4,6,8\}\pi/L=\{1.05, 1.58, 2.11\}$ GeV respectively.
%
As in Fig.~\ref{fig:quasi-soft-factor}, the results at different large $\gamma$ are consistent with each other,
demonstrating that the asymptotic limit is stable within errors.
%
We also compare the intrinsic soft function extracted from the lattice to the one-loop result in Eq.~(\ref{eq:S_ratio_1loop}),
with $\alpha_s(\mu=1/b_\perp)$ evolving from $\alpha_s(\mu=2\;{\rm GeV})\approx 0.3$. {The shaded band corresponds to the scale uncertainty  of $\alpha_s$: $\mu\in [1/\sqrt{2},\sqrt{2}]\times 1/b_\perp$. }
%
Notice that the $b_\perp$ dependence of the former comes purely from the  lattice simulation, while that for the latter is from perturbation theory. 
{For ease of comparison, we also tabulate the results for the soft function in the supplemental material~\cite{supplemental}. }



\begin{figure}[!th]
\begin{center}
\includegraphics[width=0.45\textwidth]{images/05a_phi2mod.pdf}
\includegraphics[width=0.45\textwidth]{images/05b_CS_kernel.pdf}
\caption{ Quasi-TMDWF (upper panel) and extracted Collins-Soper kernel (lower panel), as functions of $b_{\perp}$. The visible  $P^z$ dependence of the quasi-TMDWF can be primarily  understood by that from the Collins-Soper kernel, as the kernel we obtained with tree level matching is consistent with up to 3-loop perturbative calculations (at small $b_{\perp}$) with the strong coupling $\alpha_s$ at the scale $1/b_{\perp}$, and also the non-perturbative result from the  pion quasi-TMDPDF. { Results based on quenched lattice calculations, labeled as ``Hermite" and ``Bernstein"~\cite{Shanahan:2020zxr},  are also shown for comparison. Errors in the lower panel correspond to the statistical errors and the systematic errors from the non-zero imaginary part as well as  the operator mixing effects.  }  }\label{fig:quasi-TMDWF}
\end{center}
\end{figure}

We can see a clear $P^z$ dependence in the quasi-TMDWF $|\phi_\ell(0,b_{\perp},P^z,\ell)|$ normalized with $\phi_\ell(0,0,P^z,0)$, as in the upper panel of Fig.~\ref{fig:quasi-TMDWF}. This dependence is related to the CS kernel as shown in Eq.~(\ref{eq:CS_kernel_z=0}), up to   possible LaMET matching effects and power corrections of order $1/\gamma^2$. Thus we use Eq.~(\ref{eq:CS_kernel_z=0}) to extract  the kernel in   the tree level approximation, and compare the result in the lower panel of Fig.~\ref{fig:quasi-TMDWF} with that of  Ref.~\cite{Shanahan:2020zxr} and up to 3-loop perturbative ones with $\alpha_s(\mu=1/b_\perp)$. { We estimate the systematic uncertainty by combining in quadrature the  contributions from the operator mixing effects, and   from the non-vanishing  imaginary part of the quasi-TMDWF which should be {cancelled by proper treatments on higher order effects}.} For details see the supplemental materials~\cite{supplemental}, in particular {Sec. C and F}. Our result is consistent with  that of Ref.~\cite{Shanahan:2020zxr}.
%, but the results based on $P^z_1/P^z_2=3/2$ and $P^z_1/P^z_2=4/2$ have some differences at certain $b_{\perp}$ which might come from insufficient statistics at $P^z$=2.11 GeV, or potential  systematic uncertainties.
